\documentclass[11pt]{amsart}
\usepackage{amssymb}
\usepackage[margin=1.15in]{geometry}
\usepackage[hidelinks]{hyperref}
\hypersetup{pdftitle={Gabber's theorem for the etale topology}}

\newtheorem{theorem}{Theorem}[section]
\newtheorem{lemma}[theorem]{Lemma}
\newtheorem{proposition}[theorem]{Proposition}
\newtheorem{corollary}[theorem]{Corollary}
\theoremstyle{definition}
\newtheorem{definition}[theorem]{Definition}
\theoremstyle{remark}
\newtheorem{remark}[theorem]{Remark}

\newcommand{\Spec}{\operatorname{Spec}}
\newcommand{\End}{\operatorname{End}}
\newcommand{\Split}{\operatorname{Split}}
\newcommand{\cl}{\operatorname{cl}}
\newcommand{\Fib}{\operatorname{Fib}}
\newcommand{\happly}{\mathsf{happly}}
\newcommand{\ap}{\mathsf{ap}}
\newcommand{\refl}{\mathsf{refl}}
\newcommand{\close}{\Diamond}
\newcommand{\closeE}{\Diamond_{\mathrm{et}}}
\newcommand{\tr}[1]{\|#1\|}
\newcommand{\Rx}{R^\times}
\newcommand{\Mat}{M_N(R)}
\newcommand{\code}[1]{\texttt{#1}}
\newcommand{\U}{\mathcal U}

\title{Gabber's theorem for the \'etale topology, internally}
\date{November 6, 2026}

\newcount\stamphour \newcount\stampmin
\stamphour=\time \divide\stamphour by 60
\stampmin=\time \multiply\stamphour by 60 \advance\stampmin by -\stamphour \divide\stamphour by 60
\edef\stamptime{\the\year-\ifnum\month<10 0\fi\the\month-\ifnum\day<10 0\fi\the\day\space
 \ifnum\stamphour<10 0\fi\the\stamphour:\ifnum\stampmin<10 0\fi\the\stampmin}
\AddToHook{shipout/foreground}{\put(0.5\paperwidth,-\dimexpr\paperheight-0.45in\relax)%
 {\makebox[0pt]{\scriptsize\ttfamily unit/doc/\jobname.tex\quad compiled \stamptime}}}

\begin{document}
\begin{abstract}
In~\cite{h2} Gabber's theorem is proved, in gerbe form and in $H^2$ form,
in synthetic algebraic geometry with the \emph{root topology}, generated by
the roots of all monic polynomials of positive degree. Here we prove it for
the \'etale topology, generated by Zariski covers and the roots of monic
\emph{unramifiable} polynomials. The \'etale theorem is not proved again:
it is reduced to the root theorem. Two facts make the reduction work. First,
an \'etale gerbe is a root gerbe: its paths, which are \'etale sheaves and
\'etale torsors, are root sheaves by Hilbert~90. Second, an \'etale stack
gerbe is a root stack gerbe: descent for objects along a root cover follows
from Hilbert~90 applied to a product of path torsors over that cover. The
sheafification for the \'etale topology is constructed, by trees, so the
theorem has no hypothesis besides the axioms. The whole argument is checked
in Unit.
\end{abstract}
\maketitle

\section{Two topologies}

We work, as in~\cite{h2}, in univalent type theory with propositional
truncation $\tr{-}$, over a local ring $R$ satisfying the axioms of
synthetic algebraic geometry. A \emph{topology} is given by a family of
propositions $(G_j)_{j\in J}$; a type $X$ is a \emph{$J$-sheaf} if for each
$j$ the constant map $c:X\to(G_j\to X)$ is an equivalence. Sheaves are
closed under $\Pi$, $\Sigma$, identity types and retracts, and being a sheaf
is a proposition.

\emph{The root topology.} For a monic $f$ of degree $d\geq1$ let
$G_f=\tr{\{x\mid f(x)=0\}}$. The root topology is generated by all such
$G_f$; its sheafification is written $\close$ (``locally''), and is
constructed by towers of root adjunctions~\cite[Prop.~1.1]{h2}. Zariski
covers need no generator: $R$ is local, so a family $g_1,\dots,g_n$
generating $1$ has merely an invertible member.

\emph{The \'etale topology.} Let $B=R[X]/(f)$, with power basis
$1,x,\dots,x^{d-1}$, and $u=f'(x)\in B$. We say that $f$ is
\emph{unramifiable} if the $d$ coordinates of $u^d$ in the power basis
generate the unit ideal of $R$. This is the condition of~\cite{CLN}, \S1,
that the discriminant-like element $\Delta(f)=D(\delta_1,\dots,\delta_d)$
be $1$, where the $\delta_k$ are the coefficients of the characteristic
polynomial of $u$: at a prime $\mathfrak p$, all $\delta_k$ lie in
$\mathfrak p$ iff $u$ is nilpotent in $B\otimes k(\mathfrak p)$, a
$k(\mathfrak p)$-algebra of dimension $d$, iff $u^d=0$ there, iff all
coordinates of $u^d$ lie in $\mathfrak p$. The \emph{\'etale topology} is
generated by the $G_f$ with $f$ unramifiable. By~\cite{CLN}, Thm.~3.1, its
covers are the classical \'etale covers, refined by trees of localisations
and standard \'etale algebras. Every \'etale generator is a root generator,
so every root sheaf is an \'etale sheaf; the converse fails.

\begin{definition}
An \emph{\'etale tree} over a type $P$ is one of
\begin{itemize}
\item a leaf, given by an element of $P$;
\item a root node: an unramifiable $f$ and, for every root $x$ of $f$, an
\'etale tree over $P$;
\item a Zariski node: $g_1,\dots,g_n\in R$ with merely some $g_i$
invertible, and for each $i$ with $g_i$ invertible an \'etale tree over $P$.
\end{itemize}
Let $\closeE P=\tr{\text{an \'etale tree over }P}$.
\end{definition}

All data are small, so $\closeE P$ lies in the universe of $P$.

\begin{proposition}\label{p:closeE}
$\closeE P$ is a proposition and an \'etale sheaf, $P\to\closeE P$, and
$\closeE P\to Q$ for every proposition $Q$ (in any universe) which is an
\'etale sheaf and receives $P$. $\closeE$ is left exact.
\end{proposition}
\begin{proof}
A leaf gives $P\to\closeE P$. \emph{Elimination}, by induction on the tree:
at a root node, a map from the roots of $f$ to $Q$ factors through $G_f$,
and $Q$ is an \'etale sheaf; at a Zariski node, $Q$ follows from the branch
at an invertible $g_i$. \emph{Sheaf}: $\closeE P$ is a proposition, so it
suffices that $G_f\to\closeE P$ gives $\closeE P$ for $f$ unramifiable. The
roots of $f$ are the points of $\Spec B$, $B=R[X]/(f)$ finitely presented,
and at each point there is merely a tree. By $Z$-choice there merely are
$h_1,\dots,h_n\in B$ generating $1$ and, for each $i$ and each point $x$
with $h_i(x)$ invertible, a tree $t_{i,x}$. The tree with a root node $f$
followed, at the root $x$, by the Zariski node $h_1(x),\dots,h_n(x)$ with
branches $t_{i,x}$, proves $\closeE P$; the $h_i(x)$ generate $1$ at every
point. \emph{Left exactness} follows, as for any such closure, by gluing
families of sheaves over $G_f$ as products~\cite[Prop.~1.1]{h2}.
\end{proof}

\begin{remark}
The Zariski node is inside the tree, after the root node, and in general
it cannot be moved out: an \'etale tree need not be refined by a single
tower of unramifiable root adjunctions followed by one Zariski cover (an
example is $\sqrt x$ adjoined over $D(x)$, with $2$ invertible, together with
$D(1-x)$). For the root topology such a flattening exists, and the proof of
the root theorem uses it. This is why the \'etale theorem is obtained by
reduction to the root theorem.
\end{remark}

Two facts are used repeatedly, for both topologies: a sheaf proposition
that holds locally holds; and if $Q$ is a sheaf proposition with $G_j\to Q$,
then $Q$.

\section{\'Etale gerbes are root gerbes}

A \emph{band} at a point $x$ of a $1$-type $X$ is a group isomorphism
$\beta_x:(x=x)\to\Rx$. A \emph{pre-torsor} is a set $T$ with an action of
$\Rx$ which is free and transitive: for $t,t'$ in $T$ there is a unique $a$
with $at=t'$. In a $1$-type with bands, each $a=b$ is a pre-torsor, $\Rx$
acting through the band~\cite{h2}.

\begin{definition}
An \emph{\'etale gerbe} is a $1$-type $X$ with a band at every point such
that $\closeE X$, $\closeE(a=b)$ for all $a,b$, and each $a=b$ is an
\'etale sheaf. It is an \emph{\'etale stack gerbe} if moreover $X$ is an
\'etale sheaf. \emph{Root gerbes} and \emph{root stack gerbes} are defined in
the same way with $\close$ and root sheaves~\cite{h2}.
\end{definition}

\begin{lemma}[weakening]\label{l:weak}
$\closeE P\to\close P$.
\end{lemma}
\begin{proof}
$\close P$ is a root sheaf proposition, hence an \'etale sheaf, and receives
$P$; apply Proposition~\ref{p:closeE}.
\end{proof}

\begin{lemma}[Hilbert 90, \'etale]\label{l:h90}
A pre-torsor $T$ which is an \'etale sheaf and satisfies $\closeE T$ is
merely inhabited.
\end{lemma}
\begin{proof}
As for the root topology~\cite{h2}. Let $L$ be the module of maps
$f:T\to R$ with $f(at)=a^{-1}f(t)$. It is a root sheaf, since $R$ is one,
whatever $T$ is. At a point $t_0$, evaluation is an isomorphism $L\cong R$;
so $\closeE T$, weakened to $\close T$ (Lemma~\ref{l:weak}), gives that
$L$ locally has a basis of rank one, and by descent of freeness for the root
topology~\cite[Thm.~1.3]{h2} $L$ merely has a basis $f$. The type
$\{t\mid f(t)=1\}$ is a proposition: if $f(t)=f(t')=1$ and $at=t'$ then
$a^{-1}=1$. It is an \'etale sheaf ($T$ is one, and $R$ is a root sheaf).
At a point $t$ of $T$, $f(t)$ is a unit, since $f$ is a basis of $L\cong
R$, and $f(f(t)\,t)=1$. So $\closeE T$ gives a point with $f(t)=1$.
\end{proof}

\begin{lemma}\label{l:paths}
A merely inhabited pre-torsor is a root sheaf.
\end{lemma}
\begin{proof}
Being a sheaf is a proposition, so we may take a point $t_0$. Then $T$ is a
retract of $\Rx$ ($a\mapsto at_0$, and $t$ goes to the unique $a$ with
$at_0=t$), and $\Rx$ is a root sheaf.
\end{proof}

\begin{proposition}\label{p:root}
An \'etale gerbe is a root gerbe, with the same carrier and bands.
\end{proposition}
\begin{proof}
$\close X$ and $\close(a=b)$ by Lemma~\ref{l:weak}. The pre-torsor $a=b$ is
an \'etale sheaf and $\closeE(a=b)$, so it is merely inhabited
(Lemma~\ref{l:h90}), hence a root sheaf (Lemma~\ref{l:paths}).
\end{proof}

\begin{theorem}[Gabber, gerbe form, \'etale]\label{t:gerbe}
The gerbe form of Gabber's theorem~\cite[Thm.~4.1]{h2} holds for families
of \'etale gerbes: given $A$ finitely presented, a family $G$ of \'etale
gerbes over $\Spec A$ and an object of $k_*G_x$ for all $x$, merely there
are $N$, a weight-one family $V_x$ on $G_x$ with fibres of rank $N$, and
Azumaya algebras $E_x=\End(V_x)^{G_x}$ with $E_x\cong\End(V_{x,g})$ for
$g:G_x$, such that $G_x\to\Split(E_x)$ is fully faithful and preserves the
bands.
\end{theorem}
\begin{proof}
Apply the root theorem to the root gerbes of Proposition~\ref{p:root}.
\end{proof}

\begin{remark}
Each $E_x$ is \'etale-locally a matrix algebra: for $g:G_x$,
$E_x\cong\End(V_{x,g})\cong\Mat$ merely, and $\closeE G_x$.
\end{remark}

\section{\'Etale stack gerbes are root stack gerbes}

The object-descent condition does not follow from Proposition~\ref{p:root}:
an \'etale sheaf need not be a root sheaf. For the carrier of an \'etale
gerbe it does, again by Hilbert~90, now for the root topology. Fix a root
generator $P=G_f$, a proposition. We need two lemmas about $P$.

\begin{lemma}\label{l:pi}
Let $T:P\to\text{pre-torsors}$. Then $\prod_{p:P}T_p$ is a pre-torsor, with
the pointwise action.
\end{lemma}
\begin{proof}
It is a set, and the action laws hold pointwise (function
extensionality). Let $s,s'$ be in the product. For each $p$ there is a
unique $\alpha_p$ with $\alpha_p s_p=s'_p$. Since $\Rx$ is a root sheaf,
$\alpha:P\to\Rx$ is constant with a unique value $a_0$; then $a_0s=s'$. If
$bs=s'$, then $b=\alpha_p$ for each $p$ by uniqueness, so $\alpha$ is the
constant map with value $b$, and $b=a_0$. The second components are paths
in a set, hence equal.
\end{proof}

\begin{lemma}\label{l:emb}
Let $X$ be a type whose identity types are root sheaves and $c:X\to(P\to
X)$ the constant map. For $x_0:X$ the fibre
$\Fib(c\,x_0)=\sum_{x:X}(c\,x=c\,x_0)$ is contractible.
\end{lemma}
\begin{proof}
For $q:c\,x=c\,y$, let $u(q):x=y$ be the preimage of $\happly(q):P\to(x=y)$
under the constant map of $x=y$, an equivalence. Then $\ap_c(u(q))=q$:
$\happly(\ap_c r)$ is the constant family $r$ (by induction on $r$), so
$\happly(\ap_c(u(q)))=\happly(q)$, and $\happly$ is an equivalence. Hence
$(x,q)\mapsto(x,u(q))$ and $(x,r)\mapsto(x,\ap_c r)$ exhibit $\Fib(c\,x_0)$
as a retract of $\sum_x(x=x_0)$, which is contractible.
\end{proof}

\begin{theorem}[descent for objects]\label{t:stack}
The carrier of an \'etale stack gerbe is a root sheaf. Hence an \'etale stack
gerbe is a root stack gerbe, with the same carrier and bands.
\end{theorem}
\begin{proof}
Let $X$ be the carrier and fix $P=G_f$. We show that every fibre
$\Fib(\varphi)=\sum_{x:X}(c\,x=\varphi)$ of $c:X\to(P\to X)$ is
contractible; then $c$ is an equivalence.

\emph{Reduction to a point.} $\Fib(\varphi)$ is an \'etale sheaf: $X$ is
one, hence so are $P\to X$ and its identity types, and $\Sigma$ preserves
sheaves. So is the proposition ``$\Fib(\varphi)$ is contractible''
($\Sigma$, $\Pi$ and identity types again). As $\closeE X$, we may assume a
point $x_0:X$.

\emph{Hilbert 90.} Let $T=\prod_{p:P}(x_0=\varphi(p))$, a pre-torsor by
Lemma~\ref{l:pi}. Its factors are root sheaves (Proposition~\ref{p:root}),
so $T$ is one. And $\close T$: since $\close T$ is a root sheaf, it suffices
to prove it given $p:P$; then $P$ is contractible, $T\simeq(x_0=\varphi(p))$,
and $\closeE(x_0=\varphi(p))$ gives $\close(x_0=\varphi(p))$
(Lemma~\ref{l:weak}). By Hilbert~90 for the root topology~\cite[\S2]{h2}, $T$ is
merely inhabited.

\emph{Conclusion.} The goal is a proposition, so we may take $s:T$. By
function extensionality $s$ is a path $c\,x_0=\varphi$, and along it
$\Fib(\varphi)$ is $\Fib(c\,x_0)$, contractible by Lemma~\ref{l:emb}.

Finally the paths of $X$ are root sheaves and $\close X$, $\close(a=b)$ by
Proposition~\ref{p:root}.
\end{proof}

\section{Gabber's theorem in \texorpdfstring{$H^2$}{H2} form}

Recall from~\cite{h2} that $H^2(X)$ is the set truncation of
$X\to\mathrm{StackGerbe}$, root stack gerbes in the first universe, with
the contracted product as group law; $\delta(E)$ is the class of the family
of splitting gerbes $\Split(E_x)$, and $\cl(G)$ the class of $G$.

\begin{theorem}[Gabber, $H^2$ form, \'etale]\label{t:main}
Let $A$ be finitely presented and $G$ a family of \'etale stack gerbes over
$\Spec A$, and $k\geq1$.
\begin{enumerate}
\item If $k_*G_x$ has an object for every $x$, then merely there are $N$
and $E:\Spec A\to\mathrm{AZ}_N$ with $\delta(E)=\cl(G)$.
\item If $k\cdot\cl(G)=0$ in $H^2(\Spec A)$, the same conclusion holds.
\end{enumerate}
\end{theorem}
\begin{proof}
By Theorem~\ref{t:stack}, $G$ is a family of root stack gerbes; apply
\cite[Thm.~4.3]{h2} for (1) and \cite[Thm.~5.6]{h2} for (2).
\end{proof}

\begin{remark}
In the proof of (1), $G_x=\Split(E_x)$ as stack gerbes, for every $x$; so
$\Split(E_x)$ is itself an \'etale stack gerbe. The classes are compared in
$H^2$ of root stack gerbes. \'Etale stack gerbes form a subtype of root
stack gerbes, with the same identity types, since being an \'etale stack
gerbe is a proposition about a root stack gerbe. We do not prove here that
the induced map $H^2_{\mathrm{et}}\to H^2$ on set truncations is injective.
\end{remark}

\section{What is formalised in Unit}

Everything above is checked in Unit, in \code{unit/sag/twisted}:
\code{./segc sag/twisted all} gives 425 modules, 0 failed, about 11\,s.
No new axiom is added: the parameters are $R$ and the SAG axioms, with the
existing Theory parameters univalence and replacement. Both sheafifications
are constructed, not assumed. The local splitting of the remark after
Theorem~\ref{t:gerbe} and the comparison in the remark after
Theorem~\ref{t:main} are not formalised, nor is the equivalence of our
definition of unramifiable with $\Delta(f)=1$, argued in Section~1.

\smallskip
\begin{center}
\begin{tabular}{ll}
Statement & File\\\hline
unramifiable polynomials & \code{Unramifiable}\\
\'etale trees, Proposition~\ref{p:closeE} & \code{EtaleTrees}, \code{EtaleTreeDescent}, \code{EtaleTreeClosure}, \code{ClosureLex}\\
Lemmas~\ref{l:weak}--\ref{l:paths}, Proposition~\ref{p:root} & \code{EtaleGabber} (\code{weaken}, \code{H90}, \code{Root})\\
Theorem~\ref{t:gerbe} & \code{EtaleGabber.Etale}\\
Lemma~\ref{l:pi} & \code{EtaleStacks.PiTorsor}\\
Lemma~\ref{l:emb}, Theorem~\ref{t:stack} & \code{EtaleStacks.Stack}\\
Theorem~\ref{t:main} (1), (2) & \code{EtaleStacks.Etale.H2}, \code{EtaleStacks.Etale.Torsion}
\end{tabular}
\end{center}

\smallskip
The final declaration is
\[\code{EtaleStacks.Etale[R,setR,axioms].Torsion[A,fp,G,stk,n0,torsion].theorem},\]
where \code{G} is a family of \'etale gerbes, \code{stk} says that each
carrier is an \'etale sheaf, and \code{torsion} is an equality
$(n_0+1)\cdot\cl(G)=0$ in $H^2(\Spec A)$. A negative test is kept: an
\'etale sheaf is not accepted as a root sheaf without proof.

\begin{thebibliography}{9}
\bibitem{h2} \emph{Gabber's theorem in $H^2$ form, internally},
  \code{unit/doc/gabber-h2.pdf}, November 2026.
\bibitem{CLN} T. Coquand, H. Lombardi, S. Neuwirth, Azumaya algebras and
  Barr's theorem, arXiv:2306.17679.
\bibitem{Gabber} O. Gabber, Some theorems on Azumaya algebras, in \emph{The
  Brauer group}, LNM 844, Springer, 1981, 129--209.
\end{thebibliography}

\end{document}
